\subsection{网格与时间步}

问题一比较三套表面加密网格。相邻两套的全场最大差已低于 $2\times10^{-5}$，题目 $7\times5$ 表的四位小数不再随再加密而改变。将最大时间步由 $0.5\,\mathrm{s}$ 收紧到 $0.25\,\mathrm{s}$ 后，温度与水分的最大变化分别为 $2.87\times10^{-8}\,{}^\circ\mathrm{C}$ 与 $1.18\times10^{-9}\,\mathrm{kg/kg}$，论文表四位不变。问题二的粗细网格比较给出同样结论：发表表对进一步加密不敏感。因此表~\ref{tab:q1T}--表~\ref{tab:q2C} 的四位显示可以作为当前模型的稳定输出。

\subsection{扩散系数非线性}

问题一若把 $D$ 冻结为 $D(C_0)$，表面含水率在 $1800\,\mathrm{s}$ 仍会与变 $D$ 相差约 $0.04$。问题二中升温使 $D$ 增大、失水使 $D$ 减小，两者竞争后中心含水率与冻结 $D$ 相差约 $0.40$。这说明变扩散系数不是可有可无的修正，正式水分场必须保留题给 $D(C)$ 或 $D(C,T)$。

\subsection{后期环境与一维近似}

第三、四问在 $4\,\mathrm{h}$ 后采用 $50^\circ\mathrm{C}$、$0.05\,\mathrm{kg/kg}$。末小时均值与该设定几乎重合，但把最后观测点原样延续、或改用末小时均值，都会使达标时间发生有界移动。正式补入表~\ref{tab:q3C}、表~\ref{tab:q4C} 时应并列给出这一情景差，避免把单一延续写成唯一物理事实。

一维径向模型忽略端面。第一问短时可以接受；问题二 $3\,\mathrm{h}$ 热扩散长度已约 $5\,\mathrm{cm}$，更长的干燥过程中端部影响会进入内部。若后续补算轴对称 $(r,z)$ 模型，应以同一物性和两端 Robin 条件比较最大含水率到达时间，而不是在没有对照的情况下宣称一维结果已经偏于保守。
